\begin{table}[htbp]
\centering
\caption{Standardized Effect Sizes}
\label{tab:sde}
\footnotesize
\begin{adjustbox}{max width=\textwidth}
\begin{tabular}{lcccccc}
\hline\hline
Outcome & $\hat{\beta}$ & SE & SD($Y$) & SDE & SE(SDE) & Classification \\
\hline
\multicolumn{7}{l}{\textit{Panel A: Pooled}} \\
ln(Employment) & -0.0854 & 0.0357 & 1.953 & -0.0437 & 0.0183 & Small negative \\
ln(Hiring) & -0.1482 & 0.0318 & 1.787 & -0.0830 & 0.0178 & Moderate negative \\
ln(Separations) & -0.1384 & 0.0321 & 1.787 & -0.0774 & 0.0180 & Moderate negative \\
\hline
\multicolumn{7}{l}{\textit{Panel B: Heterogeneous (by EITC generosity)}} \\
ln(Emp) --- High generosity ($\geq$20\%) & -0.0079 & 0.0873 & 2.086 & -0.0038 & 0.0419 & Null \\
ln(Emp) --- Low generosity ($<$20\%) & -0.1118 & 0.0346 & 1.925 & -0.0581 & 0.0180 & Moderate negative \\
\hline\hline
\end{tabular}
\end{adjustbox}
\begin{tablenotes}[flushleft]
\small
\item \textit{Notes:} \textbf{Country:} United States. \textbf{Research question:} Do state Earned Income Tax Credit supplements increase Hispanic employment in low-wage administrative support occupations relative to non-Hispanic workers and higher-wage sectors? \textbf{Policy mechanism:} State EITC supplements provide a refundable or nonrefundable tax credit equal to 3--45\% of the federal EITC, directly increasing the after-tax return to work for low-income earners in the phase-in range and reducing the net cost of formal employment for workers near the eligibility threshold. \textbf{Outcome definition:} Log stable employment (beginning-of-quarter count from QWI) measuring workers employed at both the beginning and end of the quarter. \textbf{Treatment:} Binary indicator for whether a state has an active EITC supplement in a given year; 31 states adopted staggered over 1987--2022. \textbf{Data:} Census Bureau Quarterly Workforce Indicators (QWI) Race/Ethnicity panel, state $\times$ industry $\times$ ethnicity $\times$ year, 2000--2022, 13,752 observations. \textbf{Method:} TWFE triple-difference (state EITC $\times$ NAICS 56 $\times$ Hispanic) with state$\times$year, industry$\times$year, and state$\times$industry$\times$ethnicity fixed effects; standard errors clustered at state level. \textbf{Sample:} States with $\geq$15 years of non-missing QWI data; industries restricted to NAICS 44, 52, 54, 56, 61, 62, 72. SDE $= \hat{\beta} / \text{SD}(Y)$ where SD($Y$) is the pre-treatment standard deviation of the outcome for Hispanic workers in NAICS 56. Classification refers to magnitude, not statistical significance: Large ($|$SDE$| > 0.15$), Moderate ($0.05$--$0.15$), Small ($0.005$--$0.05$), Null ($< 0.005$). 
\end{tablenotes}
\end{table}
